\chapter{Population Dynamics: Poverty, Communalism and Sex Ratio}

\section{Poverty: Old and New Forms}

\begin{KeyConcept}
\begin{itemize}
\item \textbf{Poverty} has been defined as \textbf{absolute} (cannot meet subsistence needs) and \textbf{relative} (below a mean income) --- but these are economic. Sociologically there are also \textbf{poverty of knowledge} and \textbf{poverty of modernity} (no access to modern institutions despite finances).
\item \textbf{Social mobility vs social exclusion}: poverty can be escaped through mobility, but \textbf{social exclusion} means \textbf{no chance of mobility} --- the excluded are not part of the system (homeless, beggars, Dalits, transgenders, AIDS patients, women) --- excluded from labour market, rituals, and political representation.
\item \textbf{New forms of poverty (post-globalisation)} are based on \textbf{social exclusion and self-exclusion}. \textbf{Self-exclusion}: the working class has money but \textbf{not the courage} to enter a mall (it is 'for the rich'); professional workers feel exploited and go for \textbf{early retirement} (investing in finance rather than work); alienation (the 'gypsy' movement after the Vietnam War) leads people to opt out of the labour market, political membership and society.
\item \textbf{Relative deprivation} (Merton): after urbanisation and individuation, people lose the community's insurance and depend on \textbf{exploitative institutions} (banks, insurance, credit cards) --- so they compare themselves with others and feel deprived, feeding social/self-exclusion (and \textbf{vaccine hesitancy} --- the deprived felt the state served only the rich/connected; the state responded with \textbf{'nudge'} incentives, e.g. a bar offering free beer for a jab). Inequality has risen so much that Marx's theory of revolution seems more practical --- resentment shows in religious revivalism, ethnic violence, and could turn into identity-based civil conflict rather than class revolution.
\end{itemize}
\end{KeyConcept}

\begin{QuickRevision}
\begin{itemize}
\item Poverty
\item absolute
\item relative
\item poverty of knowledge
\item poverty of modernity
\item Social mobility vs social exclusion
\item social exclusion
\item no chance of mobility
\item New forms of poverty (post-globalisation)
\item social exclusion and self-exclusion
\item Self-exclusion
\item not the courage
\item early retirement
\item Relative deprivation
\end{itemize}
\end{QuickRevision}

\section{Communalism and Communal Violence}

\begin{KeyConcept}
\begin{itemize}
\item \textbf{Communalism} = the \textbf{fundamental distinction between groups based on religion or primordial identity}. It is a modern phenomenon (pre-British violence was ethnic/sporadic); Abrahamic religions were once thought more prone to it, but now \textbf{all religions are developing a 'political theology'} (Hindutva making uniform rules like an Abrahamic religion; Buddhist violence against Muslims in Sri Lanka/Myanmar).
\item \textbf{Communal violence has three dimensions} (CSDS): \textbf{communal riots}, \textbf{hate speech}, and \textbf{structural violence} (economic boycott --- refusing a Muslim cab-driver/delivery-boy).
\item \textbf{Reasons}: the \textbf{trauma of Partition}; colonial rule (entering a modern economy where groups blame each other for their backwardness); \textbf{relative deprivation} (the Godhra riot's economic basis --- Muslims in real estate); the use of religion and caste in politics (regional/caste parties); \textbf{history-writing} (a political act --- right-wing vs Marxist; popular history books like Bikram Sampath's Savarkar); and the \textbf{commodification of women}.
\item \textbf{Critique}: Marxists call communal violence \textbf{anti-human and anti-development} (Randhir Singh) --- people of the same (proletarian) class fight each other; \textbf{Ashis Nandy} (anti-modernist) says it is a \textbf{product of modernization} --- people lost community but did not get individualism's values, so they transform frustration into caste/religious conflict.
\item \textbf{Post-globalisation change}: from \textbf{collective punishment} (anonymous mass riots) to \textbf{individualised punishment} (mob lynching); the \textbf{culture of impunity} (posting lynching videos with the belief that one will not be convicted --- because of political support and vote-banks); anxious youth (relative deprivation, unemployment) are foot-soldiers, no longer controlled by traditional structures. \textbf{J. P. S. Uberoi}: the existence of different communities does \emph{not} cause violence --- it comes when people \textbf{define their identity solely on religion/community}.
\end{itemize}
\end{KeyConcept}

\begin{QuickRevision}
\begin{itemize}
\item Poverty: absolute/relative; poverty of knowledge/modernity; social exclusion (no mobility) + self-exclusion (opting out); relative deprivation (Merton).
\item Communalism = fundamental distinction by religion/primordial identity; modern phenomenon; 3 dimensions (riot, hate speech, structural violence).
\item Reasons: Partition trauma, colonial rule, relative deprivation (Godhra), religion-politics, history-writing, commodification of women; Marxist critique (anti-human) + Nandy (modernization).
\item Post-globalisation: collective $\rightarrow$ individualised punishment (lynching); culture of impunity; anxious youth; Uberoi (identity-definition causes violence).
\end{itemize}
\end{QuickRevision}

\section{Sex Ratio: Hypergamy, Feticide and the Dimaru Estates}

\begin{KeyConcept}
\begin{itemize}
\item India's \textbf{sex ratio} (especially the child sex ratio) is very skewed. Family planning began in the \textbf{1960s} (Rajkumari Amrit Kaur preferred 'moral' means; \textbf{sterilisation} was socially accepted; the 28-day coloured kits failed because reproductive health was taboo and child marriage common); the \textbf{1991 census} first confirmed the skewed sex ratio.
\item Sociologists found female feticide/infanticide before the census: \textbf{Louis Dumont} (hypergamy --- women of the upper strata 'missing'); \textbf{D. F. Pocock} (Patidar hypergamy); \textbf{Lata Mani} (Sati --- only $\sim$250 cases and declining, but the \emph{middle castes} were Sanskritising and adopting it, which is why reformers protested).
\item \textbf{M. N. Srinivas's 'Dimaru (daughter-eliminating) estates'}: economically prosperous states with a very skewed sex ratio --- Punjab, Haryana, Himachal, Gujarat, Maharashtra, Arunachal, Uttarakhand.
\item \textbf{Monica Dasgupta}: as the mother's \textbf{education increases, female feticide/sex-selective abortion increases} (educated mothers control fertility, so choose fewer (male) children). Reasons --- \textbf{economic} (the low value of female labour where family labour is not used; expectation of support in old age) and \textbf{cultural} (lineage carried only by the male child; the male child needed in religious rituals to avoid hell; the bride must be younger than the groom --- so a skewed ratio stays invisible).
\item \textbf{Barbara Miller's 'Bermuda Triangle of Girls'} (West-Central UP and adjoining districts): a \textbf{northern model} (low demand for female labour in agriculture) vs a \textbf{southern model} (high demand for female labour). \textbf{Amitabh Kundu \& Mahesh Sahu}: the West-Central UP districts had excess access to \textbf{amniocentesis/sonography} facilities (near Delhi) used for sex selection and illegal abortion; also the idea of \textbf{'additional sons'} (physical force to dominate), access to sex-selective technology, and higher urban living costs $\rightarrow$ sex-selective abortion; family planning itself (limiting to two children) fed the preference for the male child.
\item \textbf{Consequences}: relaxing of \textbf{gotra exogamy and village exogamy}; \textbf{import of brides} from Bihar, the North-East, Kerala (objectifying women); \textbf{reduced dowry demand}; changed inheritance (women getting more rights); more education for women; and \textbf{inter-caste marriage (by boys only)}. \textbf{North Indian marriage} = kanyadan (donation); \textbf{South Indian} = exchange of women (cross-cousin marriage, e.g. the maternal uncle) --- sapinda/incestuous marriage is banned under the Hindu Marriage Act except as custom.
\item \textbf{'Son meta-preference'}: the cultural preference for a son (to carry the lineage and perform rites) means couples keep having daughters until a son is born --- the daughters are then neglected. Relatedly, \textbf{Beteille's 'double movement of caste'}: the dominant castes pursue \textbf{Sanskritisation in society} (claiming higher ritual status) while simultaneously \textbf{de-Sanskritising before the state} (claiming OBC status for welfare benefits) --- both at the same time.
\end{itemize}
\end{KeyConcept}

\begin{QuickRevision}
\begin{itemize}
\item Family planning 1960s (sterilisation; 1991 census confirmed skewed sex ratio).
\item Dumont (hypergamy), Pocock (Patidar), Lata Mani (Sati $\sim$250, middle-caste Sanskritisation); Srinivas's Dimaru estates (Punjab/Haryana/Himachal/Gujarat/Maharashtra/Arunachal/Uttarakhand).
\item Dasgupta: education $\rightarrow$ more sex-selective abortion; economic + cultural factors; Miller's Bermuda Triangle (northern vs southern model); Kundu \& Sahu (amniocentesis, additional sons); family planning $\rightarrow$ male preference.
\item Consequences: relaxed exogamy, bride import, reduced dowry, inheritance changes, women's education, inter-caste (boys only); kanyadan vs exchange of women.
\end{itemize}
\end{QuickRevision}

\section{Reproductive Health}

\begin{KeyConcept}
\begin{itemize}
\item \textbf{Reproductive health} is measured through \textbf{ante-natal (ANC) care, institutional delivery, a skilled attendant}, and the absence of \textbf{reproductive morbidities}. India has improved in the last 20 years but remains low.
\item \textbf{Institutional delivery}: over 90\% in urban non-slum areas (paid from savings), lower in rural and lowest in slums (paid from wages --- \textbf{40\% of slum household income is spent on maternal health}). Quality differs between private and government hospitals, so the poor distrust hospitals; they also see ANC as preventive (not curative) and follow custom.
\item \textbf{Contraceptive use}: only $\sim$52\% in slum/rural areas; lower among SC/ST (lowest pre-natal care) and Muslims (61\% vs Sikhs 82\%, Hindus 72.8\%). \textbf{M. A. Hall \& R. B. Stephenson} (western India): only sterilisation is socially acceptable --- and \textbf{husbands make the decisions} about reproductive health.
\item \textbf{Women's empowerment and reproductive health}: women in \textbf{nuclear families} use institutional delivery more (than joint families with in-laws); \textbf{education} improves reproductive health (lower fertility, more ANC/contraceptives); \textbf{domestic violence} (and marital rape) during pregnancy lowers pre-natal care and leads to unintended pregnancy.
\item \textbf{Adolescent girls (10--19)}: child marriage, the discouraging of pre-marital sex, and the taboo on discussing reproductive health give them the worst reproductive health (with no rural--urban difference).
\end{itemize}
\end{KeyConcept}

\begin{QuickRevision}
\begin{itemize}
\item Reproductive health: ANC, institutional delivery, skilled attendant, morbidities; slums worst (40\% income on maternal health); private vs government quality.
\item Contraceptives $\sim$52\% (slums/rural); SC/ST lowest; Muslims lowest; Hall \& Stephenson (only sterilisation acceptable; husbands decide).
\item Nuclear family + education $\rightarrow$ better; domestic violence/marital rape $\rightarrow$ worse; adolescent girls (10--19) worst (child marriage, taboo).
\end{itemize}
\end{QuickRevision}
